\documentclass[10pt,a4paper]{article}
\usepackage[T1]{fontenc}
\usepackage[utf8]{inputenc}
\usepackage{lmodern}
\usepackage[margin=22mm,headheight=14pt]{geometry}
\usepackage{microtype}
\usepackage{amsmath,amssymb}
\usepackage{longtable,booktabs,array,calc}
\usepackage{graphicx}
\usepackage{xcolor}
\usepackage{xurl}
\usepackage{fancyhdr}
\usepackage{enumitem}
\usepackage{needspace}
\usepackage{fvextra}
\usepackage{hyperref}
\definecolor{reportblue}{HTML}{164B70}
\hypersetup{colorlinks=true,linkcolor=reportblue,urlcolor=reportblue,
  pdftitle={Qwen3.5 Inference on the SpaceMiT
K1},pdfauthor={riscv-accl project}}
\urlstyle{tt}
\setlength{\parindent}{0pt}
\setlength{\parskip}{4pt}
\setlength{\emergencystretch}{3em}
\setlength{\LTpre}{8pt}
\setlength{\LTpost}{8pt}
\setcounter{secnumdepth}{-1}
\setcounter{tocdepth}{1}
\renewcommand{\arraystretch}{1.10}
\setlist{nosep,leftmargin=1.5em}
\pagestyle{fancy}
\fancyhf{}
\fancyhead[L]{\small\textcolor{reportblue}{riscv-accl / SpaceMiT K1}}
\fancyhead[R]{\small 2026-10-14}
\fancyfoot[C]{\small\thepage}
\providecommand{\tightlist}{\setlength{\itemsep}{0pt}\setlength{\parskip}{0pt}}
\usepackage{color}
\usepackage{fancyvrb}
\newcommand{\VerbBar}{|}
\newcommand{\VERB}{\Verb[commandchars=\\\{\}]}
\DefineVerbatimEnvironment{Highlighting}{Verbatim}{commandchars=\\\{\}}
% Add ',fontsize=\small' for more characters per line
\newenvironment{Shaded}{}{}
\newcommand{\AlertTok}[1]{\textcolor[rgb]{1.00,0.00,0.00}{\textbf{#1}}}
\newcommand{\AnnotationTok}[1]{\textcolor[rgb]{0.38,0.63,0.69}{\textbf{\textit{#1}}}}
\newcommand{\AttributeTok}[1]{\textcolor[rgb]{0.49,0.56,0.16}{#1}}
\newcommand{\BaseNTok}[1]{\textcolor[rgb]{0.25,0.63,0.44}{#1}}
\newcommand{\BuiltInTok}[1]{#1}
\newcommand{\CharTok}[1]{\textcolor[rgb]{0.25,0.44,0.63}{#1}}
\newcommand{\CommentTok}[1]{\textcolor[rgb]{0.38,0.63,0.69}{\textit{#1}}}
\newcommand{\CommentVarTok}[1]{\textcolor[rgb]{0.38,0.63,0.69}{\textbf{\textit{#1}}}}
\newcommand{\ConstantTok}[1]{\textcolor[rgb]{0.53,0.00,0.00}{#1}}
\newcommand{\ControlFlowTok}[1]{\textcolor[rgb]{0.00,0.44,0.13}{\textbf{#1}}}
\newcommand{\DataTypeTok}[1]{\textcolor[rgb]{0.56,0.13,0.00}{#1}}
\newcommand{\DecValTok}[1]{\textcolor[rgb]{0.25,0.63,0.44}{#1}}
\newcommand{\DocumentationTok}[1]{\textcolor[rgb]{0.73,0.13,0.13}{\textit{#1}}}
\newcommand{\ErrorTok}[1]{\textcolor[rgb]{1.00,0.00,0.00}{\textbf{#1}}}
\newcommand{\ExtensionTok}[1]{#1}
\newcommand{\FloatTok}[1]{\textcolor[rgb]{0.25,0.63,0.44}{#1}}
\newcommand{\FunctionTok}[1]{\textcolor[rgb]{0.02,0.16,0.49}{#1}}
\newcommand{\ImportTok}[1]{#1}
\newcommand{\InformationTok}[1]{\textcolor[rgb]{0.38,0.63,0.69}{\textbf{\textit{#1}}}}
\newcommand{\KeywordTok}[1]{\textcolor[rgb]{0.00,0.44,0.13}{\textbf{#1}}}
\newcommand{\NormalTok}[1]{#1}
\newcommand{\OperatorTok}[1]{\textcolor[rgb]{0.40,0.40,0.40}{#1}}
\newcommand{\OtherTok}[1]{\textcolor[rgb]{0.00,0.44,0.13}{#1}}
\newcommand{\PreprocessorTok}[1]{\textcolor[rgb]{0.74,0.48,0.00}{#1}}
\newcommand{\RegionMarkerTok}[1]{#1}
\newcommand{\SpecialCharTok}[1]{\textcolor[rgb]{0.25,0.44,0.63}{#1}}
\newcommand{\SpecialStringTok}[1]{\textcolor[rgb]{0.73,0.40,0.53}{#1}}
\newcommand{\StringTok}[1]{\textcolor[rgb]{0.25,0.44,0.63}{#1}}
\newcommand{\VariableTok}[1]{\textcolor[rgb]{0.10,0.09,0.49}{#1}}
\newcommand{\VerbatimStringTok}[1]{\textcolor[rgb]{0.25,0.44,0.63}{#1}}
\newcommand{\WarningTok}[1]{\textcolor[rgb]{0.38,0.63,0.69}{\textbf{\textit{#1}}}}
\RecustomVerbatimEnvironment{Highlighting}{Verbatim}{breaklines=true,fontsize=\small,commandchars=\\\{\}}
\begin{document}
\begin{titlepage}
\vspace*{15mm}
{\large\textcolor{reportblue}{RISC-V LOCAL INFERENCE}}\par
\vspace{12mm}
{\Huge\bfseries Qwen3.5 Inference on the SpaceMiT K1\par}
\vspace{8mm}
{\Large Infrastructure Optimization and Performance Evaluation\par}
\vspace{12mm}
{\large Evidence cutoff: October 14, 2026 (Asia/Singapore)\par}
\vspace{14mm}
\begin{tabular}{@{}ll@{}}
\textbf{Platform} & MUSE-Pi-Pro, K1/X60, nominal 16 GB RAM \\
\textbf{Models} & Qwen3.5 2B and 4B \\
\textbf{Implementation} & SpaceMiT llama.cpp fork \\
\textbf{Package} & Technical report and evidence appendix \\
\end{tabular}
\vfill
\textcolor{reportblue}{\rule{\textwidth}{1pt}}\par
This report compares one compiled default build with the original compatibility baseline.
Long-input results have one pair per model. Formal acceptance gaps remain.\par
\vspace{6mm}
\href{https://github.com/MoyaMryia/riscv-accl}{github.com/MoyaMryia/riscv-accl}
\end{titlepage}
{\small\tableofcontents}
\clearpage
\textbf{Evidence cutoff: October 14, 2026, Asia/Singapore}

\hypertarget{abstract}{%
\section{Abstract}\label{abstract}}

The project provides one K1 build for the two tested Qwen3.5 models.
Version \nolinkurl{k1-final-20261013-r3} sets the selected parameters in
the source. Normal startup requires no performance flags or experiment
variables. The same build passes native checks, state checks, natural
code answers, and relocation checks.

A new comparison uses the original vendor commit \nolinkurl{5ad05d8}
with two compatibility changes. Both arms use identical Q4\_0 weights,
F16 cache, inputs, and output histories. At 512/2048 input tokens, 2B
decode throughput increases by 26.69\%/34.68\%. The corresponding 4B
increases are 62.17\%/73.37\%. Each short shape has three pairs. At 8192
input tokens, first-token time decreases by 51.77\% on 2B and 65.17\% on
4B. Each long shape has only one pair. These long observations do not
establish repeat stability.

The default is a balanced choice within a frozen search. It is not a
claim of a global optimum for all workloads. Model origin, full
conversion, long state correctness, and formal multi-turn quality retain
acceptance gaps. The
\href{https://github.com/MoyaMryia/riscv-accl/blob/main/spacemit/reports/snapshots/2026-10-12/SUBMISSION-REPORT.pdf}{October
12 report} preserves the previous report in full. Its source and SHA-256
receipts remain available in the same snapshot directory.

\hypertarget{requirements-and-status}{%
\section{1. Requirements and status}\label{requirements-and-status}}

\href{https://github.com/MoyaMryia/riscv-accl/blob/main/Requirements.md}{Requirements.md}
is the sole current acceptance reference. It requires both specified
model bases, local inference, the vendor fork, matched improvements, and
reproducible deployment materials. It does not set an 8 GB device
requirement or a 32k/64k context requirement. The
\href{https://github.com/MoyaMryia/riscv-accl/blob/main/spacemit/reports/2026-10-13-requirements-gap.md}{gap
report} retains the complete requirement map.

\begin{longtable}[]{@{}ll@{}}
\toprule
\begin{minipage}[b]{0.47\columnwidth}\raggedright
Status\strut
\end{minipage} & \begin{minipage}[b]{0.47\columnwidth}\raggedright
Verified scope or limit\strut
\end{minipage}\tabularnewline
\midrule
\endhead
\begin{minipage}[t]{0.47\columnwidth}\raggedright
Implemented\strut
\end{minipage} & \begin{minipage}[t]{0.47\columnwidth}\raggedright
One complete source tree, one production build, fixed defaults,
dependency checks, and a common launcher\strut
\end{minipage}\tabularnewline
\begin{minipage}[t]{0.47\columnwidth}\raggedright
Verified\strut
\end{minipage} & \begin{minipage}[t]{0.47\columnwidth}\raggedright
Both current GGUF files run locally on K1 without a cloud API\strut
\end{minipage}\tabularnewline
\begin{minipage}[t]{0.47\columnwidth}\raggedright
Verified\strut
\end{minipage} & \begin{minipage}[t]{0.47\columnwidth}\raggedright
28 matched performance requests and six complete natural code
answers\strut
\end{minipage}\tabularnewline
\begin{minipage}[t]{0.47\columnwidth}\raggedright
Verified\strut
\end{minipage} & \begin{minipage}[t]{0.47\columnwidth}\raggedright
Exact native dumps, 64 short state records, and two relocated model
launches\strut
\end{minipage}\tabularnewline
\begin{minipage}[t]{0.47\columnwidth}\raggedright
Pending\strut
\end{minipage} & \begin{minipage}[t]{0.47\columnwidth}\raggedright
Full official model origin and reproducible conversion of both final
weights\strut
\end{minipage}\tabularnewline
\begin{minipage}[t]{0.47\columnwidth}\raggedright
Pending\strut
\end{minipage} & \begin{minipage}[t]{0.47\columnwidth}\raggedright
Repeated long-context stability, long state acceptance, and the fixed
multi-turn quality suite\strut
\end{minipage}\tabularnewline
\begin{minipage}[t]{0.47\columnwidth}\raggedright
Not implemented in this release\strut
\end{minipage} & \begin{minipage}[t]{0.47\columnwidth}\raggedright
GPU acceleration, KV compression, and memory reduction through page
allocation\strut
\end{minipage}\tabularnewline
\bottomrule
\end{longtable}

A passing Unicode code task does not establish general answer quality.
Earlier shared model errors remain relevant to formal quality
acceptance. No external competition submission is recorded.

\hypertarget{platform-and-software-identity}{%
\section{2. Platform and software
identity}\label{platform-and-software-identity}}

The tested board is a MUSE-Pi-Pro with eight K1/X60 CPU cores. Its
nominal RAM is 16 GB; Linux reports approximately 15.528 GiB. Recorded
CPU frequency is 1.6 GHz. Tests use an ordinary user in tmux under the
common board lock. No device permission or system change forms part of
this campaign.

\begin{longtable}[]{@{}ll@{}}
\toprule
\begin{minipage}[b]{0.47\columnwidth}\raggedright
Item\strut
\end{minipage} & \begin{minipage}[b]{0.47\columnwidth}\raggedright
Identity\strut
\end{minipage}\tabularnewline
\midrule
\endhead
\begin{minipage}[t]{0.47\columnwidth}\raggedright
Final version\strut
\end{minipage} & \begin{minipage}[t]{0.47\columnwidth}\raggedright
\nolinkurl{k1-final-20261013-r3}\strut
\end{minipage}\tabularnewline
\begin{minipage}[t]{0.47\columnwidth}\raggedright
Vendor source\strut
\end{minipage} & \begin{minipage}[t]{0.47\columnwidth}\raggedright
\texttt{a990751d\allowbreak{}55a4c54a\allowbreak{}cf2bb77d\allowbreak{}44282c20\allowbreak{}93652359}\strut
\end{minipage}\tabularnewline
\begin{minipage}[t]{0.47\columnwidth}\raggedright
Complete source count\strut
\end{minipage} & \begin{minipage}[t]{0.47\columnwidth}\raggedright
3271 files\strut
\end{minipage}\tabularnewline
\begin{minipage}[t]{0.47\columnwidth}\raggedright
Final patch SHA-256\strut
\end{minipage} & \begin{minipage}[t]{0.47\columnwidth}\raggedright
\texttt{b03d193f\allowbreak{}618d8303\allowbreak{}ce563c4f\allowbreak{}cd4e0372\allowbreak{}6a1f6f45\allowbreak{}2dcc10ec\allowbreak{}edb1a793\allowbreak{}8550477b}\strut
\end{minipage}\tabularnewline
\begin{minipage}[t]{0.47\columnwidth}\raggedright
Server SHA-256\strut
\end{minipage} & \begin{minipage}[t]{0.47\columnwidth}\raggedright
\texttt{f254f174\allowbreak{}919c0a97\allowbreak{}ef34366f\allowbreak{}116705f6\allowbreak{}d22574ca\allowbreak{}b4c2bfc3\allowbreak{}9928c0d3\allowbreak{}15b5bfbd}\strut
\end{minipage}\tabularnewline
\begin{minipage}[t]{0.47\columnwidth}\raggedright
Policy SHA-256\strut
\end{minipage} & \begin{minipage}[t]{0.47\columnwidth}\raggedright
\texttt{aeef14e2\allowbreak{}356dc9e3\allowbreak{}3b1fc6ab\allowbreak{}595545eb\allowbreak{}67af8189\allowbreak{}4f2c7cbd\allowbreak{}28294862\allowbreak{}8453bdb0}\strut
\end{minipage}\tabularnewline
\begin{minipage}[t]{0.47\columnwidth}\raggedright
Compiler\strut
\end{minipage} & \begin{minipage}[t]{0.47\columnwidth}\raggedright
GCC 14.2.0 with vendor IME support\strut
\end{minipage}\tabularnewline
\begin{minipage}[t]{0.47\columnwidth}\raggedright
Runtime\strut
\end{minipage} & \begin{minipage}[t]{0.47\columnwidth}\raggedright
SPERT 0.6.2 and \nolinkurl{libspine_tcm.so.3.0.1}; exact file hashes
remain in the build receipt\strut
\end{minipage}\tabularnewline
\begin{minipage}[t]{0.47\columnwidth}\raggedright
Build\strut
\end{minipage} & \begin{minipage}[t]{0.47\columnwidth}\raggedright
Release, IME1 enabled, OpenMP disabled, Native disabled, ZBA enabled,
repack disabled, CPU backend\strut
\end{minipage}\tabularnewline
\bottomrule
\end{longtable}

The two model files use Q4\_0 weights and Q4\_0 embeddings. Their names
contain MTP tensors, but this release uses direct decode. The package
does not contain weights or vendor runtime libraries. Exact dependencies
do not yet have a unified public download archive.

\begin{longtable}[]{@{}ll@{}}
\toprule
\begin{minipage}[b]{0.47\columnwidth}\raggedright
Model\strut
\end{minipage} & \begin{minipage}[b]{0.47\columnwidth}\raggedright
GGUF SHA-256\strut
\end{minipage}\tabularnewline
\midrule
\endhead
\begin{minipage}[t]{0.47\columnwidth}\raggedright
\nolinkurl{Qwen3.5-2B-MTP-Q4_0-embQ4_0-dv64k.gguf}\strut
\end{minipage} & \begin{minipage}[t]{0.47\columnwidth}\raggedright
\texttt{56f0ddd9\allowbreak{}0dfa0e24\allowbreak{}56af6cb4\allowbreak{}718ece41\allowbreak{}9678d21a\allowbreak{}75d65e27\allowbreak{}353319b5\allowbreak{}e431a2fa}\strut
\end{minipage}\tabularnewline
\begin{minipage}[t]{0.47\columnwidth}\raggedright
\nolinkurl{Qwen3.5-4B-MTP-Q4_0-embQ4_0-dv64k.gguf}\strut
\end{minipage} & \begin{minipage}[t]{0.47\columnwidth}\raggedright
\texttt{0adf6cce\allowbreak{}5df53921\allowbreak{}606033c3\allowbreak{}9a13fae2\allowbreak{}054543b7\allowbreak{}8af0e7ff\allowbreak{}98860e86\allowbreak{}f078ac0f}\strut
\end{minipage}\tabularnewline
\bottomrule
\end{longtable}

Runtime fallback does not prove that inference buffers use real TCM. The
earlier sudo device audit does not establish a TCM bandwidth or
inference gain. This campaign does not add such a claim.

\hypertarget{original-baseline-and-compatibility}{%
\section{3. Original baseline and
compatibility}\label{original-baseline-and-compatibility}}

The original source is
\texttt{5ad05d87\allowbreak{}96cb094e\allowbreak{}e23a3a01\allowbreak{}52874f8a\allowbreak{}2b3060f8}
from the SpaceMiT fork. Its K1 build includes an IME2 file and fails.
The first compatibility change removes that single CMake source-list
entry. The resulting build cannot load the current reduced MTP draft
head. The second change loads \nolinkurl{d2t} and uses the stored draft
head width. Neither arm enables the draft decoder.

Only the CMake list and the tensor loader in
\nolinkurl{src/models/qwen35.cpp} change. The main output head, graph
code, and all compute kernels retain their original bytes. The
compatibility patch SHA-256 is
\texttt{2e8c40a4\allowbreak{}a0567873\allowbreak{}80e37750\allowbreak{}bab00577\allowbreak{}7580978d\allowbreak{}3d9a58d9\allowbreak{}01e911dc\allowbreak{}12fe01d7}.
Both failures and both compatibility records remain in the evidence.
This is a compatibility baseline, not an unmodified binary.

Both arms use four threads, batch 32, context 12288, one slot, and F16
K/V. They use the same compiler and runtime libraries. The measured gain
spans the original software state to the final software state. It
excludes historical weight quantization changes. It does not isolate
each individual kernel contribution.

\hypertarget{balanced-default-selection}{%
\section{4. Balanced default
selection}\label{balanced-default-selection}}

The frozen search uses one verified RC1 build and 104 requests. Four
preflight requests check both models on the original and candidate
builds. The mask stage tests eight combinations on two models and two
lengths, with two rounds. The parameter stage tests 2/4 threads and
batch sizes 16/32/64/128. Four natural answers complete the selection
checks. The parameter stage retains four threads and batch 32. No
additional parameter confirmation or fallback is required.

The balanced score gives equal weight to both models, both lengths,
TTFT, and decode time. TTFT means client time to the first output token.
A difference within 3\% of the lowest score is unresolved under the
frozen rule. This rule prefers fewer enabled candidates within that
interval. Mask 7 has the lowest observed score, 0.8904876472483813. Mask
3 scores 0.915536301962804 and lies within that interval. The rule
therefore selects mask 3. The selection score uses the same-build C0
control, not the original baseline.

\begin{longtable}[]{@{}ll@{}}
\toprule
Compiled default & Value\tabularnewline
\midrule
\endhead
Threads and batch threads & 4 / 4\tabularnewline
Batch and physical batch & 32 / 32\tabularnewline
Context and slots & 12288 / 1\tabularnewline
Cache and attention & F16 K/V, flash attention enabled\tabularnewline
GPU layers and decode & 0, direct decode\tabularnewline
Common paths & Wide RVV attention and IME route 3 enabled\tabularnewline
Candidate mask & 3: B1-r2 and C1 enabled; C2-r2 disabled\tabularnewline
\bottomrule
\end{longtable}

The launcher clears inherited experiment variables. The measured final
server receives only the model, alias, host, and port. Runtime logs and
server properties confirm the actual defaults.

\hypertarget{matched-performance-results}{%
\section{5. Matched performance
results}\label{matched-performance-results}}

Performance requests use complete fixed input token lists and 65 fixed
output tokens. A single-token bias forces identical decode histories
across versions. This synthetic workload measures performance, not
answer quality. Natural answers use a separate request without this
bias.

All requests have zero prompt-cache reuse. The tests do not clear the
operating-system file cache. Each server starts before one request;
server load time is outside request timing. All 28 requests complete and
generate exactly 65 tokens. Input hashes, output tokens, full text, and
finite timings match the protocol.

Short shapes use three pairs with alternating arm order. Long shapes use
one pair per model, with opposite order between models. Tables show
arithmetic means of per-request metrics. Each slash separates original
and final values, in that order. Time reductions and throughput
increases use different formulas.

\begin{longtable}[]{@{}lrrrrr@{}}
\toprule
Model & Input & Pairs & TTFT, s & Server prefill, s & Prefill,
tokens/s\tabularnewline
\midrule
\endhead
2B & 512 & 3 & 23.835 / 19.556 & 23.831 / 19.551 & 21.485 /
26.188\tabularnewline
2B & 2048 & 3 & 115.639 / 81.633 & 115.632 / 81.626 & 17.712 /
25.090\tabularnewline
2B & 8192 & 1 & 846.321 / 408.162 & 846.303 / 408.145 & 9.680 /
20.071\tabularnewline
4B & 512 & 3 & 63.697 / 49.865 & 63.692 / 49.860 & 8.039 /
10.269\tabularnewline
4B & 2048 & 3 & 354.441 / 213.497 & 354.434 / 213.490 & 5.778 /
9.593\tabularnewline
4B & 8192 & 1 & 3234.172 / 1126.535 & 3234.155 / 1126.518 & 2.533 /
7.272\tabularnewline
\bottomrule
\end{longtable}

\begin{longtable}[]{@{}lrrrrr@{}}
\toprule
Model & Input & Decode, tokens/s & Throughput change & Request, s & Peak
RSS, MiB\tabularnewline
\midrule
\endhead
2B & 512 & 3.387 / 4.292 & +26.69\% & 43.031 / 34.714 & 2471.97 /
2469.73\tabularnewline
2B & 2048 & 3.108 / 4.186 & +34.68\% & 136.554 / 97.169 & 2472.29 /
2469.96\tabularnewline
2B & 8192 & 2.239 / 3.600 & +60.82\% & 875.359 / 426.221 & 2473.30 /
2471.19\tabularnewline
4B & 512 & 1.315 / 2.133 & +62.17\% & 113.128 / 80.346 & 5342.20 /
5340.58\tabularnewline
4B & 2048 & 1.133 / 1.965 & +73.37\% & 411.793 / 246.579 & 5342.20 /
5340.88\tabularnewline
4B & 8192 & 0.724 / 1.485 & +104.91\% & 3323.894 / 1170.324 & 5342.88 /
5342.11\tabularnewline
\bottomrule
\end{longtable}

\begin{longtable}[]{@{}lrrrr@{}}
\toprule
Model & Input & TTFT reduction & Prefill throughput change & Decode time
reduction\tabularnewline
\midrule
\endhead
2B & 512 & 17.96\% & +21.89\% & 21.04\%\tabularnewline
2B & 2048 & 29.41\% & +41.66\% & 25.72\%\tabularnewline
2B & 8192 & 51.77\% & +107.35\% & 37.82\%\tabularnewline
4B & 512 & 21.72\% & +27.74\% & 38.34\%\tabularnewline
4B & 2048 & 39.77\% & +66.02\% & 42.32\%\tabularnewline
4B & 8192 & 65.17\% & +187.09\% & 51.20\%\tabularnewline
\bottomrule
\end{longtable}

RSS is resident memory for the server process. The table reports the
mean of sampled request peaks, with a 0.2 s sample interval. It does not
measure whole-machine memory or an instantaneous hardware peak. The
independent summary retains means, medians, minima, maxima, and each
request identity. No shape has a clear TTFT or decode-time regression
under the frozen threshold. The threshold is the greater of 3\% and the
full original-arm variation range. Three samples do not establish a
population confidence interval. One long pair does not establish
stability or complete long-context acceptance.

\hypertarget{correctness-and-failure-records}{%
\section{6. Correctness and failure
records}\label{correctness-and-failure-records}}

The final build passes 104 route, 95 attention, 432 SSM, and 158 GDN
native cases. Its dump hashes match the verified RC1 reference. Two
models produce 64 bitwise-equal full-vocabulary state records. These
records cover RS 0/3, reset states, and prefixes 32/33/64/96 at context
4096. They do not replace long state acceptance.

Six complete natural answers cover original, RC1 C0, and final arms for
both models. Each answer stops naturally and passes the frozen Unicode
task's 106 functional checks. All 2B answers have identical text and 508
output tokens. The 4B original answer has 595 tokens and different text.
The RC1 and final 4B answers have identical text and 521 tokens. Both 4B
texts pass the functional checks. This single task does not establish
broad quality or multi-turn acceptance.

The r1 delivery passes numerical checks but fails an argument-parser
GPU-default assertion. The r2 delivery fixes that assertion but starts
with four automatic slots. The r3 source fixes the server default to one
slot and receives a fresh complete build. Both failed versions retain
their source, build identities, and logs. No manual performance flag
bypasses these checks.

The previous full-combination regression remains 3.0220477\% for
4B/2048/65 against B1 under its original conditions. The new search does
not erase that failed gate. The IME scalar schedule remains excluded
after its failed performance screen. Earlier MTP output differences
require direct decode when exact output identity is required. GPU
offload and page-allocation memory benefits remain unqualified.
Predictor errors and wide intervals remain separate from measured model
results.

\hypertarget{deployment-and-evidence-package}{%
\section{7. Deployment and evidence
package}\label{deployment-and-evidence-package}}

The package contains the complete source, production binaries, fixed
launcher, policy, dependency guide, and complete evidence. Both models
use the same binary and defaults. Only the model path changes. Run this
command in tmux from the extracted version directory:

\begin{Shaded}
\begin{Highlighting}[]
\FunctionTok{bash}\NormalTok{ package/with{-}lock.sh python3 package/start.py {-}{-}model /path/to/model.gguf}
\end{Highlighting}
\end{Shaded}

Use \nolinkurl{--spert} and \nolinkurl{--tcm} only when dependency
locations differ. These options select library locations, not
performance profiles. The launcher checks binary, model, and runtime
hashes. The final build passes two model launches after binary
relocation.

The archive is
\nolinkurl{lab/final-default-20261013/release/k1-final-20261013-r3.tar.gz}.
Its external sidecar records the archive SHA-256 without a
self-reference. The evidence appendix records the final archive identity
after the archive check. The Git text view omits large dumps, thread
samples, and binaries. Strict receipt replay requires the retained
complete archive. The local check verifies 338 comparison receipts and
119 final-validation receipts. The separate search check verifies 951
search receipts and 19 native receipts.

\hypertarget{historical-evidence-and-remaining-acceptance}{%
\section{8. Historical evidence and remaining
acceptance}\label{historical-evidence-and-remaining-acceptance}}

The
\href{https://github.com/MoyaMryia/riscv-accl/blob/main/spacemit/reports/snapshots/2026-10-12/SUBMISSION-REPORT.pdf}{October
12 snapshot} retains all previous report claims and caveats. Its
\href{https://github.com/MoyaMryia/riscv-accl/blob/main/spacemit/reports/snapshots/2026-10-12/SUBMISSION-EVIDENCE.md}{evidence
appendix} and
\href{https://github.com/MoyaMryia/riscv-accl/blob/main/spacemit/reports/snapshots/2026-10-12/receipts.json}{receipts}
preserve the previous record. Historical percentages with different
controls must not be added or multiplied. The
\href{https://github.com/MoyaMryia/riscv-accl/blob/main/spacemit/reports/2026-10-07-release-verification.md}{October
7 release} and
\href{https://github.com/MoyaMryia/riscv-accl/blob/main/spacemit/reports/2026-10-13-unified-release-candidate.md}{October
13 RC1} remain comparison records. The new r3 package supplies the fixed
default delivery for this campaign.

\clearpage

\hypertarget{evidence-appendix}{%
\section{Evidence appendix}\label{evidence-appendix}}

Evidence cutoff: October 14, 2026, Asia/Singapore. This appendix
supports the
\href{https://github.com/MoyaMryia/riscv-accl/blob/main/spacemit/reports/SUBMISSION-REPORT.md}{current
report}. The
\href{https://github.com/MoyaMryia/riscv-accl/blob/main/spacemit/reports/snapshots/2026-10-12/SUBMISSION-EVIDENCE.md}{October
12 appendix} retains the previous claim map in full.

\hypertarget{current-claim-map}{%
\subsection{1. Current claim map}\label{current-claim-map}}

\begin{longtable}[]{@{}lll@{}}
\toprule
\begin{minipage}[b]{0.30\columnwidth}\raggedright
Claim\strut
\end{minipage} & \begin{minipage}[b]{0.30\columnwidth}\raggedright
Records\strut
\end{minipage} & \begin{minipage}[b]{0.30\columnwidth}\raggedright
Check and limit\strut
\end{minipage}\tabularnewline
\midrule
\endhead
\begin{minipage}[t]{0.30\columnwidth}\raggedright
One compiled default for both models\strut
\end{minipage} & \begin{minipage}[t]{0.30\columnwidth}\raggedright
\href{https://github.com/MoyaMryia/riscv-accl/blob/main/spacemit/reports/raw/k1-final-default-20261014-r3/final-r3/package/default-policy.json}{Policy},
\href{https://github.com/MoyaMryia/riscv-accl/blob/main/spacemit/reports/raw/k1-final-default-20261014-r3/final-r3/build/build-receipt.json}{build
receipt}\strut
\end{minipage} & \begin{minipage}[t]{0.30\columnwidth}\raggedright
\href{https://github.com/MoyaMryia/riscv-accl/blob/main/spacemit/bench/final-default/build_default.py}{Source
generator}; 3271 source files\strut
\end{minipage}\tabularnewline
\begin{minipage}[t]{0.30\columnwidth}\raggedright
Actual default startup\strut
\end{minipage} & \begin{minipage}[t]{0.30\columnwidth}\raggedright
\href{https://github.com/MoyaMryia/riscv-accl/blob/main/spacemit/reports/raw/k1-final-default-20261014-r3/final-validation-r3/summary.json}{Final
validation},
\href{https://github.com/MoyaMryia/riscv-accl/blob/main/spacemit/reports/raw/k1-final-default-20261014-r3/final-validation-r3/relocated-2B.server.log}{2B
log},
\href{https://github.com/MoyaMryia/riscv-accl/blob/main/spacemit/reports/raw/k1-final-default-20261014-r3/final-validation-r3/relocated-4B.server.log}{4B
log}\strut
\end{minipage} & \begin{minipage}[t]{0.30\columnwidth}\raggedright
\href{https://github.com/MoyaMryia/riscv-accl/blob/main/spacemit/bench/final-default/effective.py}{Default
checker}; no performance flags or experiment variables\strut
\end{minipage}\tabularnewline
\begin{minipage}[t]{0.30\columnwidth}\raggedright
Original compatibility boundary\strut
\end{minipage} & \begin{minipage}[t]{0.30\columnwidth}\raggedright
\href{https://github.com/MoyaMryia/riscv-accl/blob/main/spacemit/reports/raw/k1-final-default-20261014-r3/baseline-load-build/build-receipt.json}{Baseline
receipt},
\href{https://github.com/MoyaMryia/riscv-accl/blob/main/spacemit/reports/raw/k1-final-default-20261014-r3/baseline-load-compatibility.patch}{compatibility
patch}\strut
\end{minipage} & \begin{minipage}[t]{0.30\columnwidth}\raggedright
Original graph and compute kernels retain their bytes; MTP remains
disabled\strut
\end{minipage}\tabularnewline
\begin{minipage}[t]{0.30\columnwidth}\raggedright
Frozen balanced selection\strut
\end{minipage} & \begin{minipage}[t]{0.30\columnwidth}\raggedright
\href{https://github.com/MoyaMryia/riscv-accl/blob/main/spacemit/reports/raw/k1-final-default-20261014-r3/search-r2/summary.json}{Search
summary},
\href{https://github.com/MoyaMryia/riscv-accl/blob/main/spacemit/reports/raw/k1-final-default-20261014-r3/search-r2/masks-cycle-0.json}{mask
decision},
\href{https://github.com/MoyaMryia/riscv-accl/blob/main/spacemit/reports/raw/k1-final-default-20261014-r3/search-verification.json}{independent
check}\strut
\end{minipage} & \begin{minipage}[t]{0.30\columnwidth}\raggedright
104 requests; mask 3 selected under the frozen 3\% tie rule\strut
\end{minipage}\tabularnewline
\begin{minipage}[t]{0.30\columnwidth}\raggedright
Matched original-to-final performance\strut
\end{minipage} & \begin{minipage}[t]{0.30\columnwidth}\raggedright
\href{https://github.com/MoyaMryia/riscv-accl/blob/main/spacemit/reports/raw/k1-final-default-20261014-r3/comparison-r3/summary.json}{Comparison
summary},
\href{https://github.com/MoyaMryia/riscv-accl/blob/main/spacemit/reports/raw/k1-final-default-20261014-r3/final-verification.json}{independent
summary}\strut
\end{minipage} & \begin{minipage}[t]{0.30\columnwidth}\raggedright
24 short requests and four long requests; each long shape has one
pair\strut
\end{minipage}\tabularnewline
\begin{minipage}[t]{0.30\columnwidth}\raggedright
Natural answer quality\strut
\end{minipage} & \begin{minipage}[t]{0.30\columnwidth}\raggedright
\href{https://github.com/MoyaMryia/riscv-accl/blob/main/spacemit/reports/raw/k1-final-default-20261014-r3/comparison-r3/answer-2B-final/answer.json}{2B
final},
\href{https://github.com/MoyaMryia/riscv-accl/blob/main/spacemit/reports/raw/k1-final-default-20261014-r3/comparison-r3/answer-4B-final/answer.json}{4B
final}\strut
\end{minipage} & \begin{minipage}[t]{0.30\columnwidth}\raggedright
106 checks per answer; one Unicode task does not establish general
quality\strut
\end{minipage}\tabularnewline
\begin{minipage}[t]{0.30\columnwidth}\raggedright
Native and default state identity\strut
\end{minipage} & \begin{minipage}[t]{0.30\columnwidth}\raggedright
\href{https://github.com/MoyaMryia/riscv-accl/blob/main/spacemit/reports/raw/k1-final-default-20261014-r3/final-validation-r3/summary.json}{Validation
summary},
\href{https://github.com/MoyaMryia/riscv-accl/blob/main/spacemit/reports/raw/k1-final-default-20261014-r3/final-validation-r3/state.json}{state
records}\strut
\end{minipage} & \begin{minipage}[t]{0.30\columnwidth}\raggedright
789 native cases and 64 short state records; long state remains
pending\strut
\end{minipage}\tabularnewline
\begin{minipage}[t]{0.30\columnwidth}\raggedright
Preserved default failures\strut
\end{minipage} & \begin{minipage}[t]{0.30\columnwidth}\raggedright
\href{https://github.com/MoyaMryia/riscv-accl/blob/main/spacemit/reports/raw/k1-final-default-20261014-r3/r1-failure-verification.json}{r1
check},
\href{https://github.com/MoyaMryia/riscv-accl/blob/main/spacemit/reports/raw/k1-final-default-20261014-r3/r2-failure-verification.json}{r2
check}\strut
\end{minipage} & \begin{minipage}[t]{0.30\columnwidth}\raggedright
Parser assertion and automatic four-slot startup remain in the
evidence\strut
\end{minipage}\tabularnewline
\begin{minipage}[t]{0.30\columnwidth}\raggedright
Previous report preserved\strut
\end{minipage} & \begin{minipage}[t]{0.30\columnwidth}\raggedright
\href{https://github.com/MoyaMryia/riscv-accl/blob/main/spacemit/reports/snapshots/2026-10-12/receipts.json}{Snapshot
receipts}\strut
\end{minipage} & \begin{minipage}[t]{0.30\columnwidth}\raggedright
Original Markdown, TeX, PDF, template, filter, and export script retain
their bytes\strut
\end{minipage}\tabularnewline
\bottomrule
\end{longtable}

\hypertarget{complete-archive-and-compact-text-view}{%
\subsection{2. Complete archive and compact text
view}\label{complete-archive-and-compact-text-view}}

The archive path is
\nolinkurl{lab/final-default-20261013/release/k1-final-20261013-r3.tar.gz}.
The archive SHA-256 is
\texttt{d2adc88a\allowbreak{}fcb39960\allowbreak{}47b6e925\allowbreak{}d40b0bcc\allowbreak{}0054a449\allowbreak{}f9616238\allowbreak{}f698eee1\allowbreak{}c092855a}.
Its size is 1,503,683,264 bytes. The internal check verifies 5331 file
receipts, 12 symbolic links, and 28 binary hashes. The
\href{https://github.com/MoyaMryia/riscv-accl/blob/main/spacemit/final-release/archive-receipt.json}{archive
receipt} records the complete check result. The
\href{https://github.com/MoyaMryia/riscv-accl/blob/main/spacemit/reports/raw/k1-final-default-20261014-r3/ARCHIVE.md}{text
evidence guide} describes omitted large files. A receipt that names an
omitted dump cannot be fully replayed from Git alone. The complete local
evidence passes all recorded final checks.

\hypertarget{repeat-the-checks}{%
\subsection{3. Repeat the checks}\label{repeat-the-checks}}

Use the retained complete archive for strict receipt replay. Run checks
in tmux. Board builds, inference, and tests must hold the fixed common
lock. Do not modify old raw records.

\begin{Shaded}
\begin{Highlighting}[]
\ExtensionTok{python3}\NormalTok{ spacemit/bench/final{-}default/verify\_final.py }\KeywordTok{\textbackslash{}}
  \ExtensionTok{{-}{-}board}\NormalTok{ lab/final{-}default{-}20261013/board }\KeywordTok{\textbackslash{}}
  \ExtensionTok{{-}{-}delivery}\NormalTok{ lab/final{-}default{-}20261013/delivery{-}r3 }\KeywordTok{\textbackslash{}}
  \ExtensionTok{{-}{-}version}\NormalTok{ r3 {-}{-}output /tmp/k1{-}final{-}verification{-}new.json}
\end{Highlighting}
\end{Shaded}

The output file must not exist. The command above uses the complete
local collection. An extracted archive stores production binaries only
in its root \nolinkurl{build/bin/}. Create the following internal link
before a receipt replay from the extracted version directory:

\begin{Shaded}
\begin{Highlighting}[]
\FunctionTok{ln}\NormalTok{ {-}s ../../../build/bin evidence/final{-}r3/build/bin}
\ExtensionTok{python3}\NormalTok{ evidence/verification{-}tools/verify\_final.py }\KeywordTok{\textbackslash{}}
  \ExtensionTok{{-}{-}board}\NormalTok{ evidence {-}{-}delivery . {-}{-}version r3 }\KeywordTok{\textbackslash{}}
  \ExtensionTok{{-}{-}output}\NormalTok{ /tmp/k1{-}extracted{-}verification{-}new.json}
\end{Highlighting}
\end{Shaded}

The link supplies the original evidence layout without a second
production build. The checker verifies request order, identities, actual
defaults, token history, cold counts, and sampled RSS. It preserves
arithmetic means and observed ranges. The frozen selection replay uses
\nolinkurl{verify_search.py} and the original selected policy. The
current package acceptance metadata differs from the pre-test manifest
only in its acceptance status. The
\href{https://github.com/MoyaMryia/riscv-accl/blob/main/spacemit/reports/raw/k1-final-default-20261014-r3/final-package-receipt.json}{package
receipt} records that change. Source, policy, and binary bytes retain
their tested hashes.

\hypertarget{acceptance-and-historical-limits}{%
\subsection{4. Acceptance and historical
limits}\label{acceptance-and-historical-limits}}

\href{https://github.com/MoyaMryia/riscv-accl/blob/main/Requirements.md}{Requirements.md}
controls acceptance. The
\href{https://github.com/MoyaMryia/riscv-accl/blob/main/spacemit/reports/2026-10-13-requirements-gap.md}{gap
map} retains official model origin, conversion, long-state, and
multi-turn gaps. The
\href{https://github.com/MoyaMryia/riscv-accl/blob/main/spacemit/reports/2026-10-12-combined-stop-audit.md}{previous
stop audit} retains the 3.0220477\% decode regression under its original
conditions. The
\href{https://github.com/MoyaMryia/riscv-accl/blob/main/spacemit/reports/2026-10-13-k1-live-inspect-and-scalar-schedule.md}{scalar
schedule screen} retains its failed performance gate. The
\href{https://github.com/MoyaMryia/riscv-accl/blob/main/spacemit/reports/2026-10-13-k1-tcm-access.md}{TCM
audit} does not establish inference performance gains. Historical
evidence remains dated and must not be combined into an unmatched
speedup.
\end{document}
